% Môn: Vật lý
% Lớp: 11
% Chương: Chương III. Điện trường
% Bài: L11C3  Bài 17. Khái niệm điện trường
% Số câu: 20
% App đọc file này trực tiếp — sửa rồi Commit trên GitHub, bấm Đồng bộ GitHub .tex

% L11C3  Bài 17. Khái niệm điện trường

% ===== Câu 1 =====
% ID: L11C3B17-01-TL
% Mức: NB
\dangbt{Chưa có dạng}
\begin{ex}
% ID: L11C3B17-01-TL
% Mức: NB
Điện trường được tạo ra bởi điện tích, là dạng vật chất tồn tại quanh điện tích và A. tác dụng lực lên mọi vật đặt trong nó. B. tác dụng lực điện lên mọi vật đặt trong nó.C. truyền lực cho các điện tích. D. truyền tương tác giữa các điện tích. $\nguon{ly11 kntt.pdf}$
\choiceTF

\loigiai{
Chọn D.

Điện trường được tạo ra bởi điện tích, là dạng vật chất tồn tại quanh điện tích và truyền tương tác giữa các điện tích.
}
\end{ex}

% ===== Câu 2 =====
% ID: L11C3B17-01-TN
% Mức: NB
\begin{ex}
% ID: L11C3B17-01-TN
% Mức: NB
Cường độ điện trường tại một điểm đặc trưng cho điện trường tại điểm đó về $\nguon{ly11 kntt.pdf}$
\choice
{phương của vectơ cường độ điện trường.}
{chiều của vectơ cường độ điện trường.}
{\True phương diện tác dụng lực.}
{độ lớn của lực điện.}
\loigiai{
Chọn C.

Cường độ điện trường tại một điểm đặc trưng cho điện trường tại điểm đó về phương diện tác dụng lực
}
\end{ex}

% ===== Câu 3 =====
% ID: L11C3B17-02-TL
% Mức: NB
\begin{ex}
% ID: L11C3B17-02-TL
% Mức: NB
Đơn vị của cường độ điện trường là $\nguon{ly11 kntt.pdf}$
\choiceTF
{\True N.}
{\True N/m.}
{\True $\mathrm{V/m}.$}
\loigiai{
\begin{itemchoice}
\item (Đúng) N.
\item (Đúng) N/m.
\item (Đúng) $\mathrm{V/m}.$. (Vì): họn C.

Đơn vị của cường độ điện trường là V/m hoặc N/C
\end{itemchoice}
}
\end{ex}

% ===== Câu 4 =====
% ID: L11C3B17-02-TN
% Mức: NB
\begin{ex}
% ID: L11C3B17-02-TN
% Mức: NB
Một điện tích điểm $Q \lt  0$ đặt trong chân không. Cường độ điện trường do điện tích $Q$ gây ra tại một điểm $M$ cách $Q$ một khoảng r có phương là đường thẳng nối $Q$ với $M$ và $Q\nguon{ly11 kntt.pdf}$
\choice
{chiều hướng từ $M$ tới $Q$ với độ lớn bằng 4 $\pi\varepsilon$ r 2. 0 $Q$}
{chiều hướng từ $M$ ra xa khỏi $Q$ với độ lớn bằng 4 $\pi\varepsilon$ r 2. 0 - $Q$}
{\True chiều hướng từ $M$ tới $Q$ với độ lớn bằng 4 $\pi\varepsilon$ r 2. 0 - $Q$}
{chiều hướng từ $M$ ra xa khỏi $Q$ với độ lớn bằng 4 $\pi\varepsilon$ r 2. 0}
\loigiai{
Chọn C.

Một điện tích điểm Q<0 đặt trong chân không. Cường độ điện trường do điện tích $Q$ gây ra tại một điểm $M$ cách $Q$ một khoảng r có phương là đường thẳng nối $Q$ với $M$ và chiều hướng từ $M$ tới $Q$ với độ lớn bằng $\dfrac{- Q}{4\pi\varepsilon_{0}r^{2}}$.
}
\end{ex}

% ===== Câu 5 =====
% ID: L11C3B17-03-TL
% Mức: NB
\begin{bt}
% ID: L11C3B17-03-TL
% Mức: NB
Đại lượng nào dưới đây không liên quan tới cường độ điện trường của một điện tích điểm $Q$ đặt tại một điểm trong chân không? A. Khoảng cách r từ $Q$ đến điểm quan sát. B. Hằng số điện của chân không. C. Độ lớn của điện tích Q. D. Độ lớn của điện tích $Q$ đặt tại điểm quan sát. $\nguon{ly11 kntt.pdf}$
\loigiai{
Chọn D.
}
\end{bt}

% ===== Câu 6 =====
% ID: L11C3B17-03-TN
% Mức: NB
\begin{ex}
% ID: L11C3B17-03-TN
% Mức: NB
Trong chân không đặt cố định một điện tích điểm $Q = 2 \cdot  10^{-13} C$.Cường độ điện trường tại một điểm $M$ cách $Q$ một khoảng $2 \mathrm{cm}$ có giá trị bằng $\nguon{ly11 kntt.pdf}$
\choice
{$2,25 \mathrm{V/m}$.}
{\True $4,5 \mathrm{V/m}$. -4}
{2, $25.10 \mathrm{V/m}$.}
{$4,5 \cdot  10^{-4} \mathrm{V/m}.$}
\loigiai{
Chọn B.

Cường độ điện trường tại $M$ là:
}
\end{ex}

% ===== Câu 7 =====
% ID: L11C3B17-04-TL
% Mức: NB
\begin{ex}
% ID: L11C3B17-04-TL
% Mức: NB
Cường độ điện trường của Trái Đất tại điểm $M$ có giá trị bằng $120 \mathrm{V/m}$. Một electron có điện tích bằng 1,6 $\cdot$  10^{-19} $C$ và khối lượng bằng 9,1 $\cdot$  10^{-31} $\mathrm{kg}$. Chứng minh rằng, trọng lực có thể được bỏ qua so với lực điện mà Trái Đất tác dụng lên electron. Lấy $g = 9$, $8 \mathrm{m/s}$ 2. $\nguon{ly11 kntt.pdf}$
\choiceTF

\loigiai{
Lực điện tác dụng lên electron có giá trị bằng:

F_(d) = qE = 1,6$\cdot$ 10^{-19}.120 = 19,2$\cdot$ 10^{-18}$N$

Trọng lực tác dụng lên electron có giá trị bằng:

$P$ = mg = 9,1$\cdot$ 10^{-31}.9,8 = 89,18$\cdot$ 10^{-31}$N$

Từ kết quả tính được ta thấy lực điện có giá trị lớn hơn rất nhiều (hàng nghìn tỉ lần) so với trọng lực. Do đó chúng ta có thể bỏ qua trọng lực trong bài toán trên.
}
\end{ex}

% ===== Câu 8 =====
% ID: L11C3B17-04-TN
% Mức: NB
\begin{ex}
% ID: L11C3B17-04-TN
% Mức: NB
Trong chân không đặt cố định một điện tích điểm Q. Một điểm $M$ cách $Q$ một khoảng r. Tập hợp những điểm có độ lớn cường độ điện trường bằng độ lớn cường độ điện trường tại $M$ là $\nguon{ly11 kntt.pdf}$
\choice
{\True mặt cầu tâm $Q$ và đi qua M.}
{một đường tròn đi qua M.}
{một mặt phẳng đi qua M.}
{các mặt cầu đi qua M.}
\loigiai{
Chọn A.

Tập hợp những điểm có độ lớn cường độ điện trường bằng độ lớn cường độ điện trường tại Mlà mặt cầu tâm Qvà đi qua M.
}
\end{ex}

% ===== Câu 9 =====
% ID: L11C3B17-05-TL
% Mức: NB
\begin{ex}
% ID: L11C3B17-05-TL
% Mức: NB
Hãy vẽ hệ đường sức điện của điện trường xung quanh một điện tích âm đặt trong chân không và nhận xét vị trí có điện trường mạnh. $\nguon{ly11 kntt.pdf}$
\choiceTF

\loigiai{
Vẽ hệ đường sức điện của một điện tích âm Q<0 đặt trong chân không.

[Hãy vẽ hệ đường sức điện của điện trường xung quanh]

Những điểm gần điện tích $Q$ có điện trường mạnh hơn những điểm ở xa, khoảng cách càng gần thì điện trường càng mạnh.
}
\end{ex}

% ===== Câu 10 =====
% ID: L11C3B17-05-TN
% Mức: NB
\begin{ex}
% ID: L11C3B17-05-TN
% Mức: NB
Cường độ điện trường tại một điểm $M$ trong điện trường bất kì là đại lượng $\nguon{ly11 kntt.pdf}$
\choice
{\True vectơ, có phương, chiều và độ lớn phụ thuộc vào vị trí của điểm M.}
{vectơ, chỉ có độ lớn phụ thuộc vào vị trí của điểm M.}
{vô hướng, có giá trị luôn dương.}
{vô hướng, có thể có giá trị âm hoặc dương.}
\loigiai{
Chọn A.

Cường độ điện trường tại một điểm $M$ trong điện trường bất kì là đại lượng vectơ, có phương, chiều và độ lớn phụ thuộc vào vị trí của điểm M.
}
\end{ex}

% ===== Câu 11 =====
% ID: L11C3B17-06-TL
% Mức: NB
\begin{ex}
% ID: L11C3B17-06-TL
% Mức: NB
Vào một ngày đẹp trời đo đạc thực nghiệm cho thấy gân bề mặt Trái Đất ở một khu vực tại Hà Nội tôn tại điện trường theo phương thẳng đứng, hướng từ trên xuống dưới, có độ lớn cường độ điện trường không đổi trong khu vực khảo sát và bằng $114 \mathrm{V/m}$. a) Hãy vẽ hệ đường sức điện của điện trường trái đất ở khu vực đó. b) Một hạt bụi mịn có điện tích 6,4 $\cdot$  10^{-19} $C$ sẽ chịu tác dụng của lực điện có phương, chiều và độ lớn như thế nào? $\nguon{ly11 kntt.pdf}$
\choiceTF

\loigiai{
a) Vẽ hệ đường sức điện của điện trường trái đất.

[Vào một ngày đẹp trời đo đạc thực nghiệm cho thấy gần bề mặt Trái Đất]

b) Lực điện sẽ có phương thẳng đứng, chiều hướng xuống dưới theo phương và chiều của điện trường. Độ lớn của lực điện bằng: $F= qE = 6,4\cdot10^{-19}.114 = 729,6\cdot10^{-19}N$
}
\end{ex}

% ===== Câu 12 =====
% ID: L11C3B17-06-TN
% Mức: NB
\begin{ex}
% ID: L11C3B17-06-TN
% Mức: NB
Những đường sức điện của điện trường xung quanh một điện tích điểm $Q \lt  0 có dạng là\nguon{ly11 kntt.pdf}$
\choice
{những đường cong và đường thẳng có chiều đi vào điện tích Q.}
{\True những đường thẳng có chiều đi vào điện tích Q.}
{những đường cong và đường thẳng có chiều đi ra khỏi điện tích Q.}
{những đường thẳng có chiều đi ra khỏi điện tích Q.}
\loigiai{
Chọn B.

Những đường sức điện của điện trường xung quanh một điện tích điểm Q<0 có dạng là những đường thẳng có chiều đi vào điện tích Q.
}
\end{ex}

% ===== Câu 13 =====
% ID: L11C3B17-07-TL
% Mức: NB
\begin{ex}
% ID: L11C3B17-07-TL
% Mức: NB
Việc thay đổi vận tốc của electron khi qua bản lái tia như nghiên cứu ở Câu 8 có làm ảnh hưởng đến sự hiển thị tín hiệu trên màn huỳnh quang do thứ tự hiển thị tín hiệu có thể bị đảo lộn hay không? hãy giải thích. $\nguon{ly11 kntt.pdf}$
\choiceTF

\loigiai{
Cần chú ỷ rằng tín hiệu vào sẽ đặt lên hiệu điện thế của bản lái tia làm biến đổi giá trị của $U$ một cách liên tục, làm cho độ tăng của vận tốc electron khi qua điện trường cũng thay đổi liên tục.

Dựa vào kết quả của Bài III.8, chúng ta có thể giảm thiểu ảnh hưởng của sự thay đổi vận tốc khi qua bản lái tia bằng cách điều chỉnh vận tốc ban đầu v₀ tới một giá trị phù hợp để giảm tỉ lệ thay đổi của vận tốc.

Ngoài ra, ta cũng có thể điều chỉnh giảm độ lớn của hiệu điện thế $U$ đặt vào bản lái tia sẽ thu được kết quả tương tự như cách làm trên.
}
\end{ex}

% ===== Câu 14 =====
% ID: L11C3B17-07-TN
% Mức: NB
\begin{ex}
% ID: L11C3B17-07-TN
% Mức: NB
Đường sức điện cho chúng ta biết về $\nguon{ly11 kntt.pdf}$
\choice
{độ lớn của cường độ điện trường của các điểm trên đường sức điện.}
{\True phương và chiều của cường độ điện trường tại mỗi điểm trên đường sức điện.}
{độ lớn của lực điện tác dụng lên điện tích thử q.}
{độ mạnh yếu của điện trường.}
\loigiai{
Chọn B.

Đường sức điện cho chúng ta biết về phương và chiều của cường độ điện trường tại mỗi điểm trên đường sức điện.
}
\end{ex}

% ===== Câu 15 =====
% ID: L11C3B17-08-TL
% Mức: TH
\begin{ex}
% ID: L11C3B17-08-TL
% Mức: TH
Khi làm thực nghiệm xác định điện trường tại một điểm $M$ gần mặt đất, người ta dùng điện tích thử $q = 4 \cdot  10^{-16} C$ xác định được lực điện tác dụng lên điện tích q có giá trị bằng 5 $\cdot$  10^{-14} $N$, có phương thẳng đứng hướng từ trên xuống dưới. Hãy tính độ lớn cường độ điện trường tại điểm M. $\nguon{ly11kntt.pdf}$
\choiceTF

\loigiai{
Độ lớn cường độ điện trường tại điểm $M$ có giá trị bằng:

$E = \dfrac{F}{q} = \dfrac{5 \cdot 10^{-14}}{4 \cdot 10^{-16}} = 125\text{V}/\text{m.}$
}
\end{ex}

% ===== Câu 16 =====
% ID: L11C3B17-09-TL
% Mức: TH
\begin{ex}
% ID: L11C3B17-09-TL
% Mức: TH
Khi phát hiện một đám mây dông có kích thước nhỏ, một trạm quan sát thời tiết đã đo được khoảng cách từ đám mây đó đến trạm cỡ bằng $6350\,\text{m}$, người ta cũng xác định được cường độ điện trường do nó gây ra tại trạm cỡ bằng $450 \mathrm{V/m}$. Hãy ước lượng độ lớn điện tích của đám mây dông đó. Coi đám mây như một điện tích điểm. $\nguon{ly11 kntt.pdf}$
\choiceTF

\loigiai{
Điện tích $Q$ của đám mây dông có thể ước lượng theo công thức
}
\end{ex}

% ===== Câu 17 =====
% ID: L11C3B17-10-TL
% Mức: TH
\begin{ex}
% ID: L11C3B17-10-TL
% Mức: TH
Hãy vẽ hệ đường sức điện của điện trường xung quanh hệ hai điện tích âm bằng nhau và xác định những vị trí có điện trường yếu. $\nguon{ly11 kntt.pdf}$
\choiceTF

\loigiai{
Vẽ hệ đường sức điện của hệ hai điện tích âm bằng nhau Q₁ = Q₂ < 0.

[Hãy vẽ hệ đường sức điện của điện trường xung quanh hệ hai điện tích âm]

Điện trường ở vùng giữa hai điện tích là rất nhỏ, đồng thời điện trường ở cách xa hai điện tích cũng sẽ giảm dần theo khoảng cách.
}
\end{ex}

% ===== Câu 18 =====
% ID: L11C3B17-11-TL
% Mức: TH
\begin{ex}
% ID: L11C3B17-11-TL
% Mức: TH
Đặt điện tích $Q_1 = +6 \cdot  10^{-8} C$ tại điểm $A$ và điện tích $Q_2 = -2 \cdot  10^{-8} C$ tại điểm $B$ cách $A$ một khoảng bằng $3 \mathrm{cm}$. Hãy xác định những điểm mà cường độ điện trường tại đó bằng 0. $\nguon{ly11 kntt.pdf}$
\choiceTF

\loigiai{
[Đặt điện tích Q1 = +6$\cdot$ 10^-$8\,\mathrm{C}$ tại điểm $A$ và điện tích Q2 = -2$\cdot$ 10^-$8\,\mathrm{C}$]

Tại một điểm bất kì $M$ trong không gian luôn tồn tại điện trường $\overset{\rightarrow}{E_{1}}$ do điện tích Q₁ gây ra và điện trường $\overset{\rightarrow}{E_{2}}$ do điện tích Q₂ gây ra. Để cường độ điện trường tại $M$ bằng 0 thì: $\overset{\rightarrow}{E_{1}}$ = -$\overset{\rightarrow}{E_{2}}$ (Hình $17.4\,\text{G}$).

- Để $\overset{\rightarrow}{E_{1}}$, $\overset{\rightarrow}{E_{2}}$ cùng phương thì điểm $M$ phải nằm trên đường thẳng đi qua hai điểm $A$ và B.

- Để $\overset{\rightarrow}{E_{1}}$, $\overset{\rightarrow}{E_{2}}$ ngược chiều thì điểm $M$ phải nằm ngoài đoạn thẳng AB.

- Để E₁ = E₂ ⇔[Đặt điện tích Q1 = +6$\cdot$ 10^-$8\,\mathrm{C}$ tại điểm $A$ và điện tích Q2 = -2$\cdot$ 10^-$8\,\mathrm{C}$] .

Vì |Q₁|>|Q₂| nên r₁>r₂ (tức là điểm $M$ phải nằm phía ngoài điểm B).

Đặt BM = r(cm), ta có AM = 3+r (cm)

Và [Đặt điện tích Q1 = +6$\cdot$ 10^-$8\,\mathrm{C}$ tại điểm $A$ và điện tích Q2 = -2$\cdot$ 10^-$8\,\mathrm{C}$]

Giải ra ta được: $r = \dfrac{3\left( {1 + \sqrt{3}} \right)}{2}\left( \text{cm} \right)$.
}
\end{ex}

% ===== Câu 19 =====
% ID: L11C3B17-12-TL
% Mức: TH
\begin{ex}
% ID: L11C3B17-12-TL
% Mức: TH
Cho tam giác ABC vuông tại $A$ có $A$ B = $3$ \mathrm{cm} $và AC =$ 4 $\mathrm{cm}$. Tại $B$ ta đặt điện tích $Q_1 = 4$,5 \cdot  10^{-8} $C$, tại $C$, ta đặt điện tích $Q_2 = 2 \cdot  10^{-8} C$.Hãy tính độ lớn của cường độ điện trường do hai điện tích này gây ra tại A. $\nguon{ly11 kntt.pdf}$
\choiceTF

\loigiai{
[Cho tam giác ABC vuông tại $A$ có AB = $3\,\mathrm{cm}$ và AC = $4\,\mathrm{cm}$]

Sử dụng công thức [Cho tam giác ABC vuông tại $A$ có AB = $3\,\mathrm{cm}$ và AC = $4\,\mathrm{cm}$] ta tính được:

Điện trường ${\overset{\rightarrow}{E}}_{1}$ do Q₁ gây ra tại $A$ có độ lớn bằng:

$E_{1} = \dfrac{4,5\cdot 10^{-8}}{4\pi\varepsilon_{0}\left( 3\cdot 10^{-2} \right)^{2}} = 450000\text{V}/\text{m}$

Điện trường ${\overset{\rightarrow}{E}}_{2}$ do Q₂ gây ra tại $A$ có độ lớn bằng:

$E_{2} = \dfrac{2\cdot 10^{-8}}{4\pi\varepsilon_{0}\left( 4\cdot 10^{-2} \right)^{2}} = 112500\text{V}/\text{m}$

Ta thấy ${\overset{\rightarrow}{E}}_{1}$ vuông góc với ${\overset{\rightarrow}{E}}_{2}$ (Hình $17.5\,\text{G}$) nên điện trường tổng hợp ${\overset{\rightarrow}{E}}_{A}$ tính được là: $E_{A} = \sqrt{E_{1}^{2} + E_{2}^{2}} = 463849\text{V}/\text{m}$
}
\end{ex}

% ===== Câu 20 =====
% ID: L11C3B17-13-TL
% Mức: TH
\begin{ex}
% ID: L11C3B17-13-TL
% Mức: TH
Hai điểm $A$ và $B$ cách nhau $6 \mathrm{cm}$. Tại $A$, đặt điện tích $Q_1 = +8 \cdot  10^{-10} C$.Tại $B$, đặt điện tích $Q_2 = +2 \cdot  10^{-10} C$.Hãy xác định những điểm mà cường độ điện trường tại đó bằng 0. $\nguon{ly11 kntt.pdf}$
\choiceTF

\loigiai{
Do điện tích Q₁ và Q₂ cùng dấu nên vị trí cần tìm nằm giữa $A$ và B.

Để E₁ = E₂⇔[Hai điểm $A$ và $B$ cách nhau $6\,\mathrm{cm}$. Tại $A$, đặt điện tích Q1 = +8$\cdot$ 10^-$10\,\mathrm{C}$].

Đặt BM = r(cm), ta có AM = 6-r (cm)

Và [Hai điểm $A$ và $B$ cách nhau $6\,\mathrm{cm}$. Tại $A$, đặt điện tích Q1 = +8$\cdot$ 10^-$10\,\mathrm{C}$] r = $2\,\mathrm{cm}$ = BM $\Rightarrow$ AM = $4\,\mathrm{cm}$

Vậy cường độ điện trường bằng 0 tại điểm $M$ trong đoạn thẳng AB, với MA = $4\,\mathrm{cm}$ và MB = $2\,\mathrm{cm}$.
}
\end{ex}
